% =====================================================================
%  CAPÍTULO 4 — ANÁLISIS Y DISCUSIÓN DE LOS RESULTADOS
% =====================================================================
\chapter{Análisis y discusión de los resultados}
\label{cap:resultados}

Este capítulo presenta y discute los resultados del análisis, organizados en torno a las cuatro
preguntas de investigación. Se comienza por una radiografía del corpus y de su cobertura de los
estándares ESRS, que responde a la primera pregunta junto con el modelado de temas; se continúa
con el tono, el sentimiento y la especificidad del discurso, y con el índice de
\textit{greenwashing}, que sostienen la segunda y la tercera; se analizan después los factores
que determinan esas señales mediante el aparato inferencial; y se cierra con la evolución del
reporting entre los regímenes de la NFRD y la CSRD, que constituye la cuarta pregunta y, en
buena medida, el hallazgo más nítido del trabajo. A lo largo de la exposición se procura
relacionar cada resultado con el marco teórico del capítulo~\ref{cap:marco}, señalando dónde
confirma la literatura previa y dónde la matiza.

\section{Descriptivos del corpus y cobertura ESRS}
\label{sec:res-descriptivos}

La primera observación, de gran relevancia para la cuarta pregunta, surge ya del simple recuento
del texto. La tabla~\ref{tab:tokens} y la figura~\ref{fig:tokens} muestran la extensión de cada
sección del informe, medida en palabras, a lo largo de los tres ejercicios. Mientras que el
informe de gestión se mantiene estable en torno a las cuarenta mil palabras, la subsección de
sostenibilidad casi se triplica en su mediana entre 2022 y 2024, pasando de unas 6.900 a más de
17.000 palabras, lo que supone un crecimiento del orden del 147\,\%. El dato es elocuente: el
crecimiento textual no afecta al informe en su conjunto, sino que se concentra de forma
específica en la parte de sostenibilidad. Las empresas están escribiendo mucho más sobre
sostenibilidad, y lo hacen justo en el tránsito hacia la obligatoriedad de la CSRD.

\begin{table}[H]
\centering
\caption{Extensión del texto por sección y año (palabras)}
\label{tab:tokens}
\begin{tabular}{llrrr}
\toprule
\textbf{Sección} & \textbf{Año} & \textbf{n} & \textbf{Media} & \textbf{Mediana} \\
\midrule
Informe de gestión & 2022 & 194 & 38.309 & 29.822 \\
Informe de gestión & 2023 & 196 & 40.617 & 33.939 \\
Informe de gestión & 2024 & 196 & 39.736 & 31.216 \\
\textbf{Sostenibilidad} & \textbf{2022} & 194 & 11.208 & 6.894 \\
\textbf{Sostenibilidad} & \textbf{2023} & 196 & 13.201 & 9.513 \\
\textbf{Sostenibilidad} & \textbf{2024} & 196 & 20.477 & 17.041 \\
\bottomrule
\end{tabular}
\\[4pt]
{\footnotesize Fuente: elaboración propia a partir del corpus.}
\end{table}

\begin{figure}[H]
\centering
\includegraphics[width=0.82\textwidth]{5a_distribucion_tokens.png}
\caption{Distribución del tamaño de los documentos por sección y año}
\label{fig:tokens}
{\footnotesize Fuente: elaboración propia.}
\end{figure}

¿Sobre qué versa todo ese texto adicional? Para responderlo se midió la cobertura de los once
pilares del diccionario ESRS, es decir, la fracción de los términos de cada categoría que
aparece en cada informe. La tabla~\ref{tab:esrs} y la figura~\ref{fig:esrs} resumen el
resultado. Tres lecturas merecen destacarse. En primer lugar, \textit{todas} las categorías
crecen entre 2022 y 2024, con el salto más pronunciado en el último año, coherente con la
inminencia de la CSRD. En segundo lugar, el cambio climático (E1) y la plantilla propia (S1)
son, con diferencia, las categorías más cubiertas —rozan el 0,56 y el 0,53 de cobertura en
2024—, lo que concuerda con que sean los temas con mayor tradición de reporte bajo los marcos
voluntarios previos. En tercer lugar, la contaminación (E2) es sistemáticamente la categoría
menos cubierta. Conviene leer este último dato con cautela: lejos de indicar un incumplimiento
generalizado, refleja el principio de doble materialidad y un sesgo sectorial real, pues el
reporte detallado de contaminación —emisiones de óxidos, sustancias peligrosas, microplásticos—
solo es material para empresas industriales o manufactureras y carece de sentido para una
financiera o una tecnológica.

\begin{table}[H]
\centering
\caption{Cobertura media del diccionario ESRS por categoría y año}
\label{tab:esrs}
\begin{tabular}{llrrrr}
\toprule
\textbf{Cat.} & \textbf{Nombre} & \textbf{2022} & \textbf{2023} & \textbf{2024} & \textbf{$\Delta$ 22--24} \\
\midrule
E1 & Cambio climático & 0,393 & 0,444 & 0,558 & +0,165 \\
S1 & Plantilla propia & 0,406 & 0,435 & 0,528 & +0,122 \\
ESRS 2 & Divulgaciones generales & 0,296 & 0,347 & 0,464 & +0,168 \\
G1 & Conducta empresarial & 0,281 & 0,301 & 0,377 & +0,096 \\
E5 & Economía circular & 0,240 & 0,280 & 0,360 & +0,120 \\
S2 & Cadena de valor & 0,193 & 0,220 & 0,313 & +0,120 \\
E3 & Agua y recursos marinos & 0,178 & 0,207 & 0,281 & +0,103 \\
S4 & Consumidores & 0,171 & 0,198 & 0,263 & +0,092 \\
E4 & Biodiversidad & 0,147 & 0,166 & 0,241 & +0,094 \\
S3 & Comunidades & 0,153 & 0,167 & 0,225 & +0,072 \\
E2 & Contaminación & 0,112 & 0,129 & 0,195 & +0,083 \\
\bottomrule
\end{tabular}
\\[4pt]
{\footnotesize Fuente: elaboración propia. La cobertura mide la anchura del vocabulario, no la profundidad ni el cumplimiento normativo.}
\end{table}

\begin{figure}[H]
\centering
\includegraphics[width=0.82\textwidth]{5a_cobertura_esrs_anual.png}
\caption{Evolución de la cobertura ESRS por categoría y año}
\label{fig:esrs}
{\footnotesize Fuente: elaboración propia.}
\end{figure}

En conjunto, esta primera radiografía ofrece ya una respuesta provisional a la primera pregunta
de investigación —el discurso se organiza en torno al clima y a las personas, con la gobernanza
y la economía circular en un segundo plano— y anticipa con claridad la cuarta: el reporting se
expande de forma generalizada y sistemática en el umbral de la CSRD.

\section{Temas predominantes: triangulación entre LDA y BERTopic}
\label{sec:res-temas}

La cobertura por diccionario indica qué pilares ESRS aparecen, pero no descubre por sí sola los
temas que estructuran el discurso. Para ello se aplicaron las dos técnicas de modelado de temas
descritas en la metodología. El modelo LDA alcanzó su coherencia máxima con veinticinco temas,
como muestra la figura~\ref{fig:coherencia}: la curva asciende de forma clara hasta ese punto y
vuelve a descender después, lo que justifica la elección. Esa cifra —superior a la que arrojaba
una muestra más reducida en fases previas del proyecto— refleja que un corpus más amplio admite
una descomposición temática más fina sin perder interpretabilidad.

\begin{figure}[H]
\centering
\includegraphics[width=0.72\textwidth]{5b_lda_coherencia.png}
\caption{Coherencia ($C_v$) del modelo LDA según el número de temas}
\label{fig:coherencia}
{\footnotesize Fuente: elaboración propia.}
\end{figure}

Los veinticinco temas resultantes admiten una lectura muy coherente con el marco ESRS. Para
facilitar su interpretación se han agrupado en siete bloques temáticos, recogidos en la
tabla~\ref{tab:bloques}; el detalle de los veinticinco temas con sus palabras características se
ofrece en el anexo~\ref{anexo:topics}. El discurso se reparte entre un bloque climático —metas y
planes de transición, energía y emisiones—, un bloque de naturaleza —agua, contaminación,
biodiversidad y comunidades—, la economía circular, lo social referido a la propia plantilla,
lo social referido a la cadena de valor, la gobernanza y la conducta, y un bloque
específicamente ligado al nuevo marco regulatorio que reúne la doble materialidad, la Taxonomía
de la Unión y el lenguaje de los propios estándares ESRS.

\begin{table}[H]
\centering
\caption{Agrupación de los temas LDA en siete bloques}
\label{tab:bloques}
\begin{tabular}{p{4.2cm}p{9.5cm}}
\toprule
\textbf{Bloque} & \textbf{Contenido} \\
\midrule
Clima & Metas y transición, energía y eficiencia, emisiones de gases de efecto invernadero (E1) \\
Naturaleza & Agua, aire y contaminación; biodiversidad y comunidades locales (E2, E3, E4) \\
Economía circular & Residuos y materiales (E5) \\
Social --- plantilla & Formación, salud y seguridad, diversidad (S1) \\
Social --- cadena de valor & Derechos humanos, debida diligencia, proveedores (S2, S3) \\
Gobernanza y conducta & Gobierno corporativo, gestión de riesgos, ética, ciberseguridad (G1, S4) \\
Marco CSRD / Taxonomía & Doble materialidad, Taxonomía de la UE, lenguaje ESRS/GRI (ESRS 2) \\
\bottomrule
\end{tabular}
\\[4pt]
{\footnotesize Fuente: elaboración propia a partir del modelo LDA ($K=25$).}
\end{table}

El segundo modelo, BERTopic, agrupa los párrafos por significado y produce una descomposición
mucho más granular: identifica 578 temas, dejando sin clasificar en torno al 40\,\% de los
párrafos —los más genéricos o de transición—. Lejos de ser un inconveniente, esta granularidad
resulta reveladora, pues el modelo aísla como temas propios subáreas muy concretas —la gestión
del agua, la biodiversidad, la doble materialidad, la Taxonomía, la ciberseguridad, la
anticorrupción— junto a numerosos temas específicos de empresas individuales. Lo importante a
efectos de la primera pregunta es la \textit{triangulación}: los grandes temas detectados por
BERTopic, representados en la figura~\ref{fig:bertopic}, coinciden con los del modelo LDA pese a
descansar en una lógica completamente distinta. Cuando dos métodos tan diferentes identifican
los mismos asuntos, la confianza en que esos asuntos vertebran realmente el discurso aumenta de
manera considerable.

\begin{figure}[H]
\centering
\includegraphics[width=0.85\textwidth]{5b_bertopic_barchart.png}
\caption{Temas más frecuentes detectados por BERTopic}
\label{fig:bertopic}
{\footnotesize Fuente: elaboración propia.}
\end{figure}

\section{Tono, sentimiento y especificidad}
\label{sec:res-tono}

El análisis del tono y la especificidad del discurso constituye el núcleo de la tercera pregunta
y aporta evidencia decisiva sobre la cuarta. Tres modelos independientes, entrenados sobre
corpus distintos, convergen en una misma conclusión, lo que confiere a este hallazgo una
solidez notable.

El primero es el diccionario de Loughran y McDonald, cuyas siete categorías de lenguaje se
recogen en la tabla~\ref{tab:lm} y en la figura~\ref{fig:lm}. La lectura conjunta es nítida: el
lenguaje afirmativo y categórico retrocede —el tono positivo cae y la modalidad fuerte
disminuye casi un 17\,\%—, mientras que el lenguaje de incertidumbre, de cautela y de
restricción aumenta de forma consistente —la incertidumbre crece cerca de un 30\,\%—. En otras
palabras, el discurso de sostenibilidad se vuelve menos rotundo y más prudente a medida que se
acerca la obligatoriedad de la CSRD.

\begin{table}[H]
\centering
\caption{Categorías del diccionario Loughran-McDonald por año (proporción media)}
\label{tab:lm}
\begin{tabular}{lrrrrrrr}
\toprule
\textbf{Año} & \textbf{Neg.} & \textbf{Pos.} & \textbf{Incert.} & \textbf{Litig.} & \textbf{Mod. fuerte} & \textbf{Mod. débil} & \textbf{Restr.} \\
\midrule
2022 & 0,0105 & 0,0144 & 0,0094 & 0,0046 & 0,0030 & 0,0018 & 0,0049 \\
2023 & 0,0108 & 0,0144 & 0,0104 & 0,0052 & 0,0028 & 0,0020 & 0,0051 \\
2024 & 0,0126 & 0,0138 & 0,0122 & 0,0051 & 0,0025 & 0,0024 & 0,0056 \\
\bottomrule
\end{tabular}
\\[4pt]
{\footnotesize Fuente: elaboración propia.}
\end{table}

\begin{figure}[H]
\centering
\includegraphics[width=0.82\textwidth]{5c_lm_ratios.png}
\caption{Evolución de las categorías de lenguaje de Loughran-McDonald}
\label{fig:lm}
{\footnotesize Fuente: elaboración propia.}
\end{figure}

El segundo modelo, FinBERT, evalúa el tono financiero general de cada frase. Su resultado,
recogido en la tabla~\ref{tab:finbert} y en la figura~\ref{fig:finbert}, apunta en la misma
dirección con una claridad incluso mayor: el tono neto —la diferencia entre frases positivas y
negativas— cae de forma monótona y sustancial, desde 0,189 en 2022 hasta 0,140 en 2024, lo que
representa una reducción de más del 25\,\%. La caída se concentra, significativamente, en el
salto de 2023 a 2024, justo en el umbral de la CSRD. El aumento del peso de las frases neutras
sugiere, además, un desplazamiento hacia un texto más descriptivo y técnico, en detrimento del
discurso evaluativo y promocional.

\begin{table}[H]
\centering
\caption{Tono financiero (FinBERT) por año}
\label{tab:finbert}
\begin{tabular}{lrrrr}
\toprule
\textbf{Año} & \textbf{\% positivo} & \textbf{\% negativo} & \textbf{\% neutral} & \textbf{Tono (pos.\,$-$\,neg.)} \\
\midrule
2022 & 23,50 & 4,65 & 71,84 & 0,1885 \\
2023 & 23,33 & 4,76 & 71,91 & 0,1857 \\
2024 & 19,60 & 5,55 & 74,85 & 0,1405 \\
\bottomrule
\end{tabular}
\\[4pt]
{\footnotesize Fuente: elaboración propia.}
\end{table}

\begin{figure}[H]
\centering
\includegraphics[width=0.82\textwidth]{5c_finbert_sentiment.png}
\caption{Evolución del tono financiero (FinBERT) por año}
\label{fig:finbert}
{\footnotesize Fuente: elaboración propia.}
\end{figure}

El tercer modelo, ClimateBERT, se aplica solo a las frases climáticas —en torno a dos de cada
cinco del corpus— y resulta el más revelador para la tercera pregunta. La tabla~\ref{tab:climate}
recoge sus tres dimensiones. Las tres se mueven en la misma dirección entre 2022 y 2024: el
discurso de \textit{oportunidad} climática cae casi cinco puntos mientras el de \textit{riesgo}
sube en proporción similar —hasta el punto de que en 2024 el riesgo ya supera a la oportunidad,
cuando en 2022 esta casi lo duplicaba—; las frases con un \textit{compromiso} concreto retroceden
del 32,8 al 26,1\,\%; y las frases \textit{específicas}, las que aportan cifras, fechas o metas,
también disminuyen. Esta convergencia es importante porque matiza la hipótesis clásica del
\textit{greenwashing}. No se observa un discurso cada vez más optimista y vacío —el patrón que
cabría esperar de un \textit{cheap talk} en aumento—, sino algo más sutil: un discurso que se
vuelve simultáneamente menos optimista, menos comprometido y menos específico. La lectura más
plausible, que el índice de la sección siguiente permite confirmar, es que el enorme crecimiento
del volumen de texto se debe en buena parte a contenido descriptivo de cumplimiento normativo
—definiciones, metodología de la doble materialidad, tablas de indicadores— más que a nuevos
compromisos cuantificados.

\begin{table}[H]
\centering
\caption{Sentimiento, compromiso y especificidad climática (ClimateBERT) por año (\%)}
\label{tab:climate}
\begin{tabular}{lrrr@{\hspace{0.8cm}}rr@{\hspace{0.8cm}}rr}
\toprule
 & \multicolumn{3}{c}{\textbf{Sentimiento}} & \multicolumn{2}{c}{\textbf{Compromiso}} & \multicolumn{2}{c}{\textbf{Especificidad}} \\
\cmidrule(lr){2-4}\cmidrule(lr){5-6}\cmidrule(lr){7-8}
\textbf{Año} & Neutral & Oport. & Riesgo & Sí & No & Sí & No \\
\midrule
2022 & 68,3 & 20,7 & 11,0 & 32,8 & 67,2 & 26,2 & 73,8 \\
2023 & 67,4 & 20,0 & 12,6 & 31,1 & 68,9 & 27,1 & 72,9 \\
2024 & 67,3 & 15,9 & 16,8 & 26,1 & 73,9 & 24,1 & 75,9 \\
\bottomrule
\end{tabular}
\\[4pt]
{\footnotesize Fuente: elaboración propia.}
\end{table}

Como tercera vía de caracterización temática, independiente del diccionario y del modelado de
temas, el clasificador FinBERT-ESG-9 confirma la imagen anterior: el cambio climático domina y
crece ligeramente, lo social y la gobernanza ocupan el segundo escalón, y las categorías de
naturaleza y contaminación permanecen estables y bajas, en plena coherencia con la cobertura por
diccionario. Resulta destacable que la categoría \textit{no-ESG} aumente su peso, lo que
respalda la hipótesis de que parte del crecimiento de la subsección de sostenibilidad bajo la
CSRD es contenido normativo y estructural que ningún clasificador asocia a una temática de
sostenibilidad concreta. La figura~\ref{fig:esg9} ilustra esta distribución.

\begin{figure}[H]
\centering
\includegraphics[width=0.82\textwidth]{5c_esg9_distribucion.png}
\caption{Distribución temática de las frases según FinBERT-ESG-9}
\label{fig:esg9}
{\footnotesize Fuente: elaboración propia.}
\end{figure}

\section{El índice de \textit{greenwashing}}
\label{sec:res-gw}

Todas las señales anteriores confluyen en el índice de \textit{greenwashing} (\gwindex{}), que
sintetiza en una sola medida la cautela del lenguaje, la vaguedad de las promesas, la escasez de
cifras y la falta de especificidad climática. Su evolución, recogida en la tabla~\ref{tab:gw} y
en la figura~\ref{fig:gw}, dibuja una trayectoria inequívoca: el índice se mantiene en valores
negativos en 2022 y 2023 y da un salto pronunciado al terreno positivo en 2024. La
descomposición por componentes, representada en la figura~\ref{fig:gwcomp}, identifica los
motores de ese deterioro: el lenguaje evasivo aumenta, la especificidad climática disminuye y,
de manera especialmente reveladora, el contenido cuantitativo \textit{también} se reduce en
2024. Este último dato desmonta una objeción evidente: podría pensarse que un informe más
extenso será necesariamente más detallado, pero ocurre lo contrario: pese a que el texto casi se
triplica, la proporción de frases con cifras concretas cae. Conviene subrayar, en sentido
inverso, que las promesas vagas sin cuantificar permanecen prácticamente estables, de modo que
el deterioro no proviene de un aumento de las promesas huecas, sino de una pérdida general de
sustancia —menos cifras, menos especificidad— acompañada de más cautela.

\begin{table}[H]
\centering
\caption{Índice de \textit{greenwashing} y sus componentes por año}
\label{tab:gw}
\begin{tabular}{lrrrrr}
\toprule
\textbf{Año} & \textbf{\gwindex{}} & \textbf{Hedging} & \textbf{R.\ cuantit.} & \textbf{Promesa vaga} & \textbf{Espec.\ clim.} \\
\midrule
2022 & $-$0,2026 & 0,0112 & 0,3234 & 0,1241 & 0,2621 \\
2023 & $-$0,1677 & 0,0124 & 0,3235 & 0,1238 & 0,2707 \\
2024 & \textbf{+0,3705} & 0,0146 & 0,3049 & 0,1226 & 0,2413 \\
\bottomrule
\end{tabular}
\\[4pt]
{\footnotesize Fuente: elaboración propia.}
\end{table}

\begin{figure}[H]
\centering
\includegraphics[width=0.78\textwidth]{5d_gwindex_evolucion.png}
\caption{Evolución del índice de \textit{greenwashing} (2022--2024)}
\label{fig:gw}
{\footnotesize Fuente: elaboración propia.}
\end{figure}

\begin{figure}[H]
\centering
\includegraphics[width=0.82\textwidth]{5d_componentes.png}
\caption{Componentes del índice de \textit{greenwashing} por año}
\label{fig:gwcomp}
{\footnotesize Fuente: elaboración propia.}
\end{figure}

El índice se relaciona, como cabía esperar, con el discurso de riesgo climático: la correlación
entre ambos es positiva y de magnitud moderada ($r=0{,}46$), mientras que con la oportunidad y
el compromiso es negativa. La correlación transversal con el tono financiero general es, en
cambio, casi nula, un punto sobre el que se volverá al discutir la tercera pregunta. La lectura
temporal es la verdaderamente relevante: el tono cae al mismo tiempo que el índice sube, de modo
que el discurso se vuelve a la vez menos optimista y más evasivo. No es el patrón clásico de
«más optimismo sin sustancia», sino el de «menos sustancia en general» bajo el nuevo marco
regulatorio.

\section{Determinantes por sector, país y tamaño}
\label{sec:res-determinantes}

Las secciones anteriores describen tendencias agregadas; esta examina su \textit{estructura}, es
decir, cómo se distribuyen las señales según las características de la empresa, lo que da
respuesta a la segunda pregunta y a buena parte de la tercera. La prueba de Kruskal-Wallis,
resumida en la tabla~\ref{tab:kruskal}, confirma que tanto el índice de \textit{greenwashing}
como el tono y la especificidad difieren de forma estadísticamente muy significativa según el
sector y según el país: todos los contrastes son significativos al uno por mil. Existe, por
tanto, una heterogeneidad real que las medias agregadas ocultaban.

\begin{table}[H]
\centering
\caption{Prueba de Kruskal-Wallis por sector y por país}
\label{tab:kruskal}
\begin{tabular}{lrrrr}
\toprule
\textbf{Variable} & \multicolumn{2}{c}{\textbf{Por sector}} & \multicolumn{2}{c}{\textbf{Por país}} \\
\cmidrule(lr){2-3}\cmidrule(lr){4-5}
 & $H$ & $\eta^2$ & $H$ & $\eta^2$ \\
\midrule
\gwindex{} & 57,39 & 0,100 & 85,26 & 0,148 \\
Tono (FinBERT) & 36,45 & 0,062 & 62,13 & 0,106 \\
Especificidad climática & 66,29 & 0,115 & 44,51 & 0,077 \\
\bottomrule
\end{tabular}
\\[4pt]
{\footnotesize Fuente: elaboración propia. Todos los contrastes con $p<0{,}001$.}
\end{table}

El detalle por sector, representado en la figura~\ref{fig:gwsector}, revela un patrón coherente
y muy interpretable. Los sectores financiero y tecnológico exhiben los valores más altos del
índice de \textit{greenwashing} y, simultáneamente, los niveles más bajos de especificidad
climática; en el extremo opuesto, las utilities, el sector inmobiliario y los materiales básicos
presentan los índices más bajos y la mayor especificidad. La lógica es transparente: las
empresas cuya huella climática directa es menor —o cuya materialidad climática es más financiera
que operativa— tienden a un discurso climático más cualitativo y genérico, mientras que las de
actividad física intensiva reportan con mayor concreción. El resultado conecta directamente con
el \textit{cherry-picking} descrito por \textcite{bingler2022}: quienes menos tienen que
contar en términos operativos son quienes más recurren al lenguaje vago.

\begin{figure}[H]
\centering
\includegraphics[width=0.85\textwidth]{5e_gwindex_supersector.png}
\caption{Índice de \textit{greenwashing} por supersector}
\label{fig:gwsector}
{\footnotesize Fuente: elaboración propia.}
\end{figure}

La dimensión geográfica, recogida en la figura~\ref{fig:gwregion}, muestra que las empresas del
Reino Unido e Irlanda y las nórdicas presentan índices más elevados que las de Europa central y,
sobre todo, que las del sur. Este patrón debe leerse con prudencia, pues algunos países cuentan
con muy pocas empresas en la muestra —de ahí que en el análisis de regresión se agrupen en
regiones—, pero la dirección general es consistente.

\begin{figure}[H]
\centering
\includegraphics[width=0.85\textwidth]{5e_gwindex_region.png}
\caption{Índice de \textit{greenwashing} por región}
\label{fig:gwregion}
{\footnotesize Fuente: elaboración propia.}
\end{figure}

El análisis de regresión permite aislar el efecto de cada factor manteniendo los demás
constantes. La primera regresión explica el índice de \textit{greenwashing}; sus coeficientes
más relevantes se recogen en la tabla~\ref{tab:reg1} y el detalle completo, en el
anexo~\ref{anexo:regresiones}. Confirma que los efectos sectoriales y regionales son robustos:
el sector financiero y el tecnológico mantienen índices significativamente superiores al de
referencia, y las regiones nórdica y británico-irlandesa, también. Dos resultados merecen
comentario. El primero es que el tamaño de la empresa, medido por su capitalización, no muestra
ningún efecto significativo sobre el índice —un resultado que matiza la intuición de que las
mayores compañías comunican necesariamente mejor—. El segundo, de gran importancia para la
cuarta pregunta, es que el año 2024 conserva un efecto positivo y significativo sobre el índice
\textit{aun después de controlar por sector, tamaño y región}: el deterioro de 2024 no es, por
tanto, un mero artefacto de la composición de la muestra, sino que apunta a un efecto genuino
asociado a la entrada en vigor de la CSRD.

\begin{table}[H]
\centering
\caption{Regresión 1: determinantes del índice de \textit{greenwashing} (coeficientes seleccionados)}
\label{tab:reg1}
\begin{tabular}{lrr}
\toprule
\textbf{Predictor} & \textbf{Coef.} & \textbf{$p$} \\
\midrule
Sector financiero & +2,034 & $<$0,001 \\
Sector tecnológico & +1,338 & 0,002 \\
Región nórdica & +0,612 & 0,022 \\
Región Reino Unido e Irlanda & +0,751 & 0,003 \\
Año 2024 & +0,510 & 0,030 \\
Capitalización (log) & $-$0,026 & 0,723 \\
\bottomrule
\end{tabular}
\\[4pt]
{\footnotesize Fuente: elaboración propia. $R^2\approx0{,}14$; categorías de referencia: sector materiales básicos, año 2022, región centro.}
\end{table}

La segunda regresión aborda de frente la hipótesis del \textit{greenwashing} clásico: ¿se asocia
un tono más optimista a un discurso menos específico? El resultado, recogido en la
tabla~\ref{tab:reg2}, es claro y contrario a esa hipótesis simple. La especificidad climática
tiene un efecto \textit{positivo} y significativo sobre el tono ($+0{,}194$): las empresas que
hablan de clima de forma más concreta tienden también a un tono financiero más positivo,
probablemente porque la concreción suele acompañar a la comunicación de logros y avances
tangibles. La hipótesis ingenua de que el optimismo encubre vaguedad no se sostiene en el corte
transversal. Lo que sí se confirma, una vez más, es el efecto temporal: el año 2024 deprime el
tono de forma robusta e independiente de cuánto y cómo hable la empresa de clima. El sector
inmobiliario y el energético aparecen, además, como los de tono más contenido.

\begin{table}[H]
\centering
\caption{Regresión 2: determinantes del tono financiero (coeficientes seleccionados)}
\label{tab:reg2}
\begin{tabular}{lrr}
\toprule
\textbf{Predictor} & \textbf{Coef.} & \textbf{$p$} \\
\midrule
Especificidad climática & +0,194 & 0,003 \\
Año 2024 & $-$0,048 & $<$0,001 \\
Región Reino Unido e Irlanda & +0,064 & $<$0,001 \\
Sector inmobiliario & $-$0,069 & $<$0,001 \\
Sector energético & $-$0,075 & 0,021 \\
Capitalización (log) & +0,000 & 0,938 \\
\bottomrule
\end{tabular}
\\[4pt]
{\footnotesize Fuente: elaboración propia. $R^2\approx0{,}16$; categorías de referencia: sector materiales básicos, año 2022, región centro.}
\end{table}

\section{Evolución del reporting: de la NFRD a la CSRD}
\label{sec:res-evolucion}

La cuarta pregunta —cómo evoluciona el reporting entre los dos regímenes normativos— ha ido
recibiendo respuesta a lo largo de todo el capítulo, pero conviene cerrarla con el contraste
formal que le da solidez estadística. La prueba pareada compara cada empresa consigo misma en
2022 y en 2024, sobre las 194 compañías presentes en ambos ejercicios, y sus resultados se
recogen en la tabla~\ref{tab:pareado} y en la figura~\ref{fig:pareado}. El veredicto es
contundente: el reporting de 2024 es significativamente más extenso, menos optimista, con más
discurso de riesgo climático y con un índice de \textit{greenwashing} más elevado que el de
2022. Cinco de las seis variables alcanzan significación al uno por mil o muy próxima a ese
nivel; solo la especificidad climática, que disminuye en la dirección esperada, no llega a ser
estadísticamente significativa, lo que la confirma como la señal más sutil del conjunto.

\begin{table}[H]
\centering
\caption{Contraste pareado 2022 frente a 2024 (194 empresas comunes)}
\label{tab:pareado}
\begin{tabular}{lrrrr}
\toprule
\textbf{Variable} & \textbf{Media 2022} & \textbf{Media 2024} & \textbf{$\Delta$} & \textbf{$p$ (Wilcoxon)} \\
\midrule
\gwindex{} & $-$0,208 & +0,339 & +0,547 & 0,005 \\
Tono (FinBERT) & 0,1885 & 0,1422 & $-$0,046 & $<$0,001 \\
Especificidad climática & 0,263 & 0,244 & $-$0,019 & 0,126 \\
Riesgo climático & 0,1106 & 0,1651 & +0,055 & $<$0,001 \\
Oportunidad climática & 0,2069 & 0,1604 & $-$0,046 & $<$0,001 \\
Palabras (sostenibilidad) & 11.208 & 20.603 & +9.395 & $<$0,001 \\
\bottomrule
\end{tabular}
\\[4pt]
{\footnotesize Fuente: elaboración propia.}
\end{table}

\begin{figure}[H]
\centering
\includegraphics[width=0.85\textwidth]{5e_pareado_2022_2024.png}
\caption{Cambios pareados entre 2022 y 2024 en las variables clave}
\label{fig:pareado}
{\footnotesize Fuente: elaboración propia.}
\end{figure}

A este contraste se suma la evidencia más visual de toda la investigación, procedente del
modelado de temas. Al seguir la frecuencia de cada tema a lo largo del tiempo, destaca con
nitidez el de la \textbf{doble materialidad}, que multiplica su presencia por más de nueve entre
2022 y 2024, como muestra la figura~\ref{fig:overtime}. La doble materialidad es el concepto
nuclear que introduce la CSRD y que no existía bajo la NFRD, de modo que su irrupción repentina
en 2024 constituye la huella textual más fuerte y cuantificable del cambio de régimen. Otros
temas estrechamente ligados al nuevo marco crecen también con fuerza —la Taxonomía de la Unión,
el riesgo climático físico, el agua y la biodiversidad—, lo que confirma que la transición no
solo aumenta el volumen del reporting, sino que reorienta su contenido hacia el lenguaje propio
de los estándares europeos.

\begin{figure}[H]
\centering
\includegraphics[width=0.85\textwidth]{5b_topics_over_time.png}
\caption{Evolución temporal de temas seleccionados (2022--2024)}
\label{fig:overtime}
{\footnotesize Fuente: elaboración propia.}
\end{figure}

En síntesis, los resultados permiten responder a las cuatro preguntas con un grado de
convergencia notable entre métodos. El discurso de sostenibilidad se organiza en torno al clima
y a las personas (RQ1); difiere de forma significativa por sector y país, con el sector
financiero y el tecnológico a la cabeza del lenguaje vago (RQ2); no responde a la hipótesis
ingenua de que el optimismo encubre vaguedad —la relación entre especificidad y tono es
positiva—, aunque sí revela una pérdida general de sustancia (RQ3); y experimenta, en el tránsito
de la NFRD a la CSRD, una transformación profunda: más texto, menos optimismo, más cautela, más
riesgo y menos concreción, con la doble materialidad como emblema del nuevo régimen (RQ4). El
sentido último de esta transformación —si es un avance hacia un reporting más honesto sobre la
incertidumbre o una señal de \textit{cheap talk} creciente— se discute en las conclusiones.

A modo de cierre del capítulo conviene mencionar la herramienta con la que se han explorado
todos estos resultados. El conjunto de hallazgos se ha integrado en un cuadro de mando
interactivo de elaboración propia que permite navegar por el corpus sin necesidad de manejar los
datos en bruto: ofrece una visión agregada del material, un explorador que muestra el perfil de
cualquiera de las 196 empresas, un mapa de los temas detectados, un comparador entre compañías y
una síntesis de los resultados por pregunta de investigación. Más allá de su utilidad como
soporte de la defensa del trabajo, este cuadro de mando ilustra el potencial de las técnicas
empleadas como instrumento de cribado: permitiría a un analista, un inversor o un supervisor
localizar con rapidez los informes cuyo perfil textual —mucho lenguaje vago, poca
especificidad— aconseja una lectura más atenta.
